% =====================================================================
%  ClassMind 开题报告 - 主文件
%  使用 xelatex 编译
% =====================================================================
\documentclass[a4paper,12pt,oneside]{ctexart}

% ----- 页面 -----
\usepackage[a4paper,top=2.7cm,bottom=2.6cm,left=2.8cm,right=2.6cm,headheight=16pt,footskip=1.2cm]{geometry}
\usepackage{fancyhdr}
\usepackage{setspace}

% ----- 字体 -----
\usepackage{fontspec}
\setmainfont{Times New Roman}
\setsansfont{Arial}
\setmonofont{Consolas}[Scale=MatchLowercase]
\setCJKmainfont{SimSun}[BoldFont=SimHei,ItalicFont=KaiTi,BoldItalicFont=SimHei]
\setCJKsansfont{SimHei}
\setCJKmonofont{NSimSun}

% ----- 颜色 / 链接 / 代码 -----
\usepackage{xcolor}
\definecolor{cmblue}{HTML}{0F4C81}
\definecolor{cmgreen}{HTML}{2E8B57}
\definecolor{cmgray}{HTML}{6B7280}
\definecolor{cmaccent}{HTML}{B8860B}
\usepackage[hidelinks]{hyperref}
\usepackage{listings}
\usepackage{algorithm}
\usepackage{algpseudocode}
\usepackage{booktabs}
\usepackage{tabularx}
\usepackage{longtable}
\usepackage{multirow}
\usepackage{array}
\usepackage{graphicx}
\usepackage{tikz}
\usetikzlibrary{shapes.geometric,arrows.meta,positioning,calc}
\usepackage{pgfplots}
\pgfplotsset{compat=1.18}
\usepackage{caption}
\usepackage{subcaption}
\usepackage{enumitem}
\usepackage{amsmath,amssymb}
\usepackage{siunitx}
\usepackage{rotating}
\usepackage{url}
\usepackage{microtype}
\usepackage{placeins}      % \FloatBarrier：强制当前位置完成所有浮动
\usepackage{float}         % 允许 [H] 强制就地

% ----- 行距 -----
\setstretch{1.45}
\setlength{\emergencystretch}{3em}

% ----- 代码样式 -----
\lstdefinestyle{cmplain}{
  basicstyle=\ttfamily\footnotesize,
  breaklines=true,
  frame=single,
  rulecolor=\color{cmgray!30},
  backgroundcolor=\color{gray!5},
  numbers=left,
  numberstyle=\tiny\color{cmgray},
  numbersep=8pt,
  showstringspaces=false,
  keepspaces=true,
  columns=fullflexible,
  xleftmargin=1.5em,
  xrightmargin=0.5em,
}
\lstdefinestyle{cmjson}{
  basicstyle=\ttfamily\footnotesize,
  commentstyle=\color{cmgray}\itshape,
  breaklines=true,
  frame=single,
  rulecolor=\color{cmgray!30},
  backgroundcolor=\color{gray!5},
  numbers=left,
  numberstyle=\tiny\color{cmgray},
  numbersep=8pt,
  showstringspaces=false,
}

% ----- 页眉页脚 -----
\pagestyle{fancy}
\fancyhf{}
\fancyhead[L]{\small 高途 ClassMind 智能排课决策系统}
\fancyhead[R]{\small 开题报告}
\fancyfoot[C]{\small -- \thepage\ --}
\renewcommand{\headrulewidth}{0.4pt}
\renewcommand{\footrulewidth}{0pt}

% ----- 章节编号：使用规范的阿拉伯数字层级 -----
\renewcommand{\thesection}{\arabic{section}}
\renewcommand{\thesubsection}{\thesection.\arabic{subsection}}
\renewcommand{\thesubsubsection}{\thesubsection.\arabic{subsubsection}}
\setcounter{secnumdepth}{3}
\setcounter{tocdepth}{2}

% ----- 标题样式：只由 ctex 管理，避免重复定义 -----
\ctexset{
  section={format=\heiti\zihao{-3}\raggedright,beforeskip=18bp,afterskip=10bp},
  subsection={format=\heiti\zihao{4}\raggedright,beforeskip=12bp,afterskip=6bp},
  subsubsection={format=\heiti\zihao{-4}\raggedright,beforeskip=8bp,afterskip=4bp}
}

% ----- 目录：仅显示两级，统一缩进、引导点与页码 -----
\makeatletter
\renewcommand*{\l@section}{\@dottedtocline{1}{0em}{2.8em}}
\renewcommand*{\l@subsection}{\@dottedtocline{2}{2.8em}{3.6em}}
\renewcommand*{\@dotsep}{1.5}
\makeatother

% ----- 自定义命令 -----
\newcommand{\code}[1]{\texttt{#1}}
\newcommand{\cm}[1]{\textcolor{cmblue}{\textbf{#1}}}
\newcolumntype{Y}{>{\raggedright\arraybackslash}X}
\newcolumntype{C}[1]{>{\centering\arraybackslash}p{#1}}
\newcolumntype{L}[1]{>{\raggedright\arraybackslash}p{#1}}

\setlength{\parindent}{2em}
\setlength{\parskip}{0pt}
\setlist{nosep,leftmargin=2em,topsep=0.4em}

% ----- caption 字体 -----
\captionsetup{font={small,rm},labelfont=bf,labelsep=quad,format=hang,skip=6pt}
\renewcommand{\arraystretch}{1.18}
\setlength{\tabcolsep}{5pt}

% ----- tikz 风格 -----
\tikzstyle{cmblock}=[draw=black,rectangle,rounded corners=1pt,minimum height=0.9cm,minimum width=2.2cm,align=center,fill=white,text=black]
\tikzstyle{cmarrow}=[-{Latex[length=2mm]},thick,draw=black]
\tikzstyle{cmnode}=[circle,draw=black,fill=white,minimum size=0.6cm,inner sep=1pt]

% =====================================================================
\begin{document}
% =====================================================================

% ----- 自动在每章/每节前清空浮动体，避免图表被推到下章 -----
\let\oldsection\section
\let\oldsubsection\subsection
\let\oldsubsubsection\subsubsection
\renewcommand{\section}{\FloatBarrier\oldsection}
\renewcommand{\subsection}{\FloatBarrier\oldsubsection}
\renewcommand{\subsubsection}{\FloatBarrier\oldsubsubsection}

% 封面页（不计入页码）
% 封面页
\begin{titlepage}
\vspace*{0.4cm}
\begin{center}
{\fontsize{26pt}{32pt}\selectfont\bfseries 高途 ClassMind\\智能排课决策系统\par}
\vspace{0.7cm}
{\LARGE\bfseries 开\quad 题\quad 报\quad 告\par}
\vspace{0.9cm}
{\Large\bfseries 基于约束规划的智能排课系统设计与实现\par}
\vspace{0.3cm}
{\Large ——以高途一对一辅导业务为应用场景\par}
\vspace{1.2cm}
\end{center}

\begin{center}
\renewcommand{\arraystretch}{1.35}
\begin{tabular}{>{\raggedleft\arraybackslash}p{3.3cm}p{9.0cm}}
{\large 项目名称：} & {\large\bfseries 高途排课脑 ClassMind} \\[0.35cm]
{\large 申报方向：} & {\large 智能体应用 / 教育科技} \\[0.35cm]
{\large 项目阶段：} & {\large MVP 第一阶段（已完成） + 第二阶段规划} \\[0.35cm]
{\large 起止时间：} & {\large 2026 年 07 月 -- 2026 年 12 月} \\[0.35cm]
{\large 团队成员：} & {\large 李帅锋（队长 / 算法 / 全栈）} \\[0.35cm]
{\large 指导教师：} & {\large ＿＿＿＿＿＿ 教授} \\[0.35cm]
{\large 参赛赛事：} & {\large 飞书高校应用创新大赛} \\
\end{tabular}
\end{center}

\vfill
\begin{center}
{\large 2026 年 7 月\par}
\end{center}
\end{titlepage}

% 承诺书页
\thispagestyle{empty}
\vspace*{2cm}
\begin{center}
{\LARGE\bfseries 原创性与使用授权声明\par}
\end{center}
\vspace{1.5cm}

\noindent 本项目（高途排课脑 ClassMind）由参赛团队独立设计、研发与撰写报告。
项目涉及的算法、源代码、演示数据与文档均为团队原创工作，不存在抄袭、
剽窃等学术不端行为。引用的公开文献、第三方库（Google OR-Tools 等）
已按学术规范在文末列出。

本项目计划参加飞书高校应用创新大赛，参赛作品如获得奖项，参赛团队
同意主办单位对项目方案、源代码、文档与演示视频进行非商业性的
展示、汇编与公开宣传。

\vspace{1.0cm}
\begin{flushright}
\begin{tabular}{l}
团队负责人（签字）：李帅锋 \\
\rule{4cm}{0.4pt}\\[0.3cm]
\qquad\quad 年\qquad 月\qquad 日
\end{tabular}
\end{flushright}
\newpage


% 前置部分：摘要、目录（不设置图目录和表目录）
\pagestyle{fancy}
\pagenumbering{Roman}
% 中文摘要与英文摘要分别成页
\clearpage
\phantomsection
\addcontentsline{toc}{section}{摘要}
\begin{center}
  {\heiti\zihao{2} 摘\quad 要}
\end{center}
\vspace{1.2em}

个性化辅导排课需要同时协调教师资质、班级可用时间、教室容量、设备条件
及临时变更等多类约束。依赖电子表格和人工核对的传统方式难以在业务变化
频繁时稳定保证无冲突，也不易解释不同排课方案之间的取舍。为此，本课题
拟研究并实现一套面向 K12 一对一及小班辅导场景的智能排课决策系统
ClassMind。

课题以约束规划和 OR-Tools CP-SAT 为核心，将“课次完整安排、教师/教室/
班级同一时段不冲突、教师资质匹配、教室容量及设备满足”等条件建模为硬
约束，将学生偏好、教师偏好、资源利用、日负载均衡和调课变更数建模为可
配置软目标。系统拟采用“业务数据层—约束求解层—独立校验层—应用交互层”
四层架构，形成候选方案生成、评分解释、异常诊断和最小影响调课的闭环，
并为第二阶段接入飞书多维表格、消息卡片和审批流程预留接口。

目前已完成可运行原型，包括数据模型、CP-SAT 求解器、本地 REST API、
四个功能页面和自动化测试。后续研究将围绕真实规模数据下的求解性能、
多目标权重敏感性、局部重排稳定性、结果可解释性和飞书场景集成展开。
预期成果包括一套可复现实验的智能排课原型、约束与评价指标体系、飞书
集成方案以及完整的技术文档。

\vspace{1em}
\noindent\textbf{关键词：}智能排课；约束规划；CP-SAT；多目标优化；
最小影响调课；可解释决策；飞书应用

\clearpage
\phantomsection
\addcontentsline{toc}{section}{Abstract}
\begin{center}
  {\bfseries\zihao{2} Abstract}
\end{center}
\vspace{1.2em}

Personalized tutoring timetabling must coordinate teacher qualifications,
class availability, room capacity, equipment requirements, and frequent
schedule changes. Spreadsheet-based workflows rely heavily on manual checks,
making it difficult to guarantee conflict-free schedules or explain the
trade-offs among alternatives. This project proposes ClassMind, an intelligent
timetabling decision system for one-to-one and small-group K12 tutoring.

The project adopts constraint programming with OR-Tools CP-SAT. Lesson
completion, teacher-room-class exclusivity, qualification matching, capacity,
and equipment requirements are modeled as hard constraints. Student and
teacher preferences, resource utilization, workload balance, and schedule
changes are modeled as configurable soft objectives. The proposed architecture
contains four layers: business data, constraint solving, independent
validation, and application interaction. It supports candidate generation,
score-based explanation, infeasibility diagnosis, and minimum-impact
rescheduling, while reserving interfaces for Feishu Base, message cards, and
approval workflows.

A runnable prototype has been completed, including the data model, CP-SAT
solver, local REST APIs, four functional pages, and automated tests. Subsequent
research will evaluate scalability, weight sensitivity, rescheduling stability,
interpretability, and Feishu integration. Expected deliverables include a
reproducible scheduling prototype, a constraint and evaluation framework, a
Feishu integration plan, and complete technical documentation.

\vspace{1em}
\noindent\textbf{Keywords:} Intelligent Timetabling; Constraint Programming;
CP-SAT; Multi-objective Optimization; Minimum-impact Rescheduling;
Explainable Decision; Feishu Application
\clearpage

% 目录页
\clearpage
\thispagestyle{fancy}
\begingroup
\renewcommand{\contentsname}{目\quad 录}
\hypersetup{linkcolor=black}
\tableofcontents
\endgroup
\clearpage


% 正文
\pagenumbering{arabic}
% ==================== 第 1 章 选题背景与研究意义 ====================
\clearpage
\section{选题背景与研究意义}

\subsection{问题背景}

个性化辅导排课并不是简单地把课程填入日历，而是一个多资源、多约束、
动态变化的组合优化问题。每个课次至少涉及课程、班级、教师、教室和时间槽
五类对象；对象之间又存在资质、容量、设备、不可用时间和偏好时段等关系。
当课次数量增加或教师临时请假时，局部调整可能触发连锁变更，人工逐项核对
容易遗漏冲突，也难以稳定复现同一套决策规则。

现有教务流程常以电子表格维护基础数据，再通过即时通信工具确认变更。该
方式适合小规模、低频变化场景，但存在三方面局限：一是约束分散在人员经验
和多个表格中，难以形成统一规则；二是方案质量依赖人工搜索顺序，无法保证
硬约束始终满足；三是方案调整缺少结构化差异记录，解释、审批和通知成本较高。

因此，本课题选择“约束求解 + 独立校验 + 可解释交互”的技术路线，把不可
违反的业务规则与可权衡的运营偏好分开建模，以机器求解负责组合搜索，以
独立校验器复核结果，以可视化界面呈现方案差异。

\subsection{研究意义}

\subsubsection{方法意义}

课题将教师、教室、班级和时间槽冲突统一表示为离散约束，并把偏好命中、
资源利用、负载均衡和变更数量纳入同一多目标框架。与只输出单一课表相比，
这种建模方式能够明确区分“必须满足”和“尽量满足”，便于分析不同权重对
方案的影响，也便于在业务规则变化时扩展模型。

\subsubsection{应用意义}

系统可将重复核对工作转化为数据校验和自动求解，帮助教务人员把精力集中到
规则确认、异常处理和方案审批。独立校验器能够对输出课表再次检查，降低因
程序实现或接口传输造成的冲突风险；局部重排机制则能在教师请假等情况下
优先保留原方案，减少通知范围。

\subsubsection{平台意义}

飞书多维表格适合维护结构化业务对象，消息卡片和审批流程适合承接方案确认
与通知。本课题先完成可独立运行的排课核心，再设计飞书集成接口，可避免把
关键求解逻辑绑定在单一前端，同时为后续形成“数据维护—智能排课—审批发布—
变更通知”的协同闭环奠定基础。

\subsection{应用场景与研究边界}

本课题聚焦 K12 一对一及小班辅导的周排课与临时调课。研究对象包括教师、
教室、班级、课程请求和离散时间槽；输出为满足全部硬约束的课表及其评分、
解释和变更明细。第一阶段不讨论学费结算、招生 CRM、在线授课、考勤和复杂
权限体系，也不把自然语言模型作为排课正确性的来源。自然语言只作为未来的
输入辅助，最终约束仍需转化为可校验的结构化数据。

% ==================== 第 2 章 国内外研究现状 ====================
\clearpage
\section{国内外研究现状}

\subsection{排课与组合优化研究}

排课问题属于典型的组合优化问题。早期研究常将课程与时段冲突转化为图着色
或整数规划模型；随后，遗传算法、模拟退火、禁忌搜索等元启发式方法被用于
较大规模实例。此类方法能够在限定时间内搜索较优方案，但通常需要针对业务
设计修复算子，若处理不当，可能生成违反硬约束的中间解。

约束规划强调以变量、定义域和约束直接描述问题。现代 CP-SAT 求解器进一步
融合约束传播、SAT 冲突学习和整数优化技术，适合表达“恰好一次”“至多一次”
和资源互斥等离散规则。其优势在于：当求解器返回可行状态时，模型中声明的
硬约束均被满足；同时可通过目标函数优化偏好与资源利用。

\subsection{智能排课系统研究}

现有排课系统大体可分为三类。第一类是面向学校固定课表的专用系统，重点
处理课程、教师和教室的周期性安排；第二类是培训机构 SaaS，强调教务流程、
客户管理与通知；第三类是面向研究的算法原型，重点比较求解策略和标准数据集。

对个性化辅导场景而言，固定课表方法难以覆盖频繁调课，通用 SaaS 的内部
求解逻辑通常不透明，而研究原型又往往缺少可直接使用的业务界面。除此之外，
“求出一张表”并不等于完成教务决策：实际使用还要求解释偏好取舍、展示受
影响对象、保留历史版本并支持审批发布。

\subsection{现有研究的不足}

结合本课题应用场景，现有方法主要存在以下不足：

\begin{enumerate}
  \item 硬约束和软偏好混合表达，导致结果是否可执行不够清晰；
  \item 侧重首次排课，对教师请假后的最小影响重排研究不足；
  \item 评价指标只给出单一总分，难以解释学生、教师与资源之间的取舍；
  \item 算法原型与协同平台脱节，数据维护、方案审批和消息通知难以形成闭环；
  \item 缺少独立于求解器的结果复核，系统正确性过度依赖单一实现。
\end{enumerate}

\subsection{本课题的切入点}

本课题不追求发明新的通用求解器，而是研究如何把成熟的 CP-SAT 能力转化为
可解释、可校验、可集成的教务决策系统。具体切入点包括：建立面向个性化
辅导的结构化约束模型；设计多策略权重与分项评分；以变更数量为主要目标
实现局部重排；用独立校验器形成双重保障；最后把求解能力封装为 API，为
飞书多维表格、消息卡片和审批流程提供稳定接口。

% ==================== 第 3 章 研究目标与主要内容 ====================
\clearpage
\section{研究目标、主要内容与关键问题}

\subsection{总体目标}

本课题拟构建一套零硬冲突、可解释、可复核、可扩展的智能排课决策系统。
系统接收结构化教师、教室、班级、课程和时间槽数据，输出满足全部硬约束的
课表；在存在多个可行方案时，根据不同业务策略优化软目标，并提供评分、
原因说明和变更明细。最终形成可独立运行的原型，同时完成与飞书协同流程的
接口设计和验证方案。

\subsection{主要研究内容}

\begin{enumerate}
  \item \textbf{业务对象与约束建模。}明确五类核心对象及其关系，把资质、
  容量、设备、不可用时段和资源互斥转化为可计算约束。
  \item \textbf{多目标排课模型。}在硬约束之外，引入学生偏好、教师偏好、
  容量适配、时间紧凑和负载均衡等软目标，并设计可配置策略。
  \item \textbf{最小影响调课。}以原课表为基线，把变更课次数量纳入目标，
  生成新课表及逐项差异。
  \item \textbf{校验与解释。}建立独立校验器，复核课次完整性和资源冲突；
  设计分项指标，使教务人员能够理解方案取舍。
  \item \textbf{应用与平台集成。}实现本地 API 和交互界面，设计飞书多维
  表格数据映射、消息卡片和审批回调方案。
  \item \textbf{实验与评价。}围绕正确性、性能、稳定性、偏好敏感性和
  可用性开展实验，形成可复现的评价流程。
\end{enumerate}

\subsection{拟解决的关键问题}

\subsubsection{硬约束的完备表达}

需要保证所有实际不可违反的规则都进入模型，并避免同一规则在前端、后端和
数据表中以不同口径重复实现。课题将以统一数据模型为入口，以求解模型和
独立校验器分别实现约束检查，降低漏约束风险。

\subsubsection{多目标之间的可解释权衡}

学生偏好、教师偏好和资源效率可能相互冲突。课题不直接比较不同权重体系下
的 CP-SAT 原始目标值，而采用统一口径的分项指标和归一化评分比较方案；
同时通过权重敏感性实验分析方案变化。

\subsubsection{动态变化下的方案稳定性}

调课不应简单地重新生成完全不同的课表。系统需要识别受影响课次，以原方案
为基线，在满足新增硬约束的前提下最小化变更数量，并输出变更原因和影响对象。

\subsection{技术指标与验收方式}

表 \ref{tab:targets} 中的数值是本课题拟达到的验收目标，不是外部行业统计。

\begin{table}[htbp]
  \centering
  \caption{研究目标与验收方式}
  \label{tab:targets}
  \small
  \begin{tabularx}{\textwidth}{L{2.6cm} L{3.2cm} Y Y}
    \toprule
    \textbf{维度} & \textbf{目标} & \textbf{验收方法} & \textbf{证据形式} \\
    \midrule
    正确性 & 可行方案硬冲突数为 0 & 求解后运行独立校验器，并覆盖不可行场景 & 自动化测试与校验报告 \\
    完整性 & 所有课程请求按规定课次排入 & 对课程请求与输出课次逐项核对 & 课次清单与缺失检查 \\
    性能 & 典型规模在设定时限内返回可行结果 & 分层构造不同规模数据并重复测量 & 运行环境、实例规模与耗时记录 \\
    稳定性 & 调课时优先减少相对基线的变更 & 设置教师不可用事件并统计变更课次 & 调课前后差异清单 \\
    可解释性 & 展示至少三类分项指标及方案差异 & 检查评分卡、解释文本和权重配置 & 界面截图与指标定义 \\
    可集成性 & 提供结构化接口并完成飞书映射设计 & 验证 JSON 输入输出及回调流程 & API 文档与集成原型 \\
    \bottomrule
  \end{tabularx}
\end{table}

% ==================== 第 4 章 研究方案与技术路线 ====================
\clearpage
\section{研究方案与技术路线}

\subsection{总体研究思路}

本课题采用“结构化输入—约束建模—候选求解—独立校验—评分解释—协同发布”
的研究思路。业务规则首先被划分为硬约束和软目标；CP-SAT 负责搜索可行解
并优化指定目标；独立校验器不复用求解器内部状态，而是直接检查输出课表；
应用层再把课表、评分和变更明细呈现给教务人员。该设计把算法正确性、业务
解释和平台交互分离，便于独立测试和迭代。

\subsection{系统总体架构}

图 \ref{fig:architecture-new} 采用黑白分层结构。外部协同平台不直接控制
求解器，而是通过应用服务层交换结构化数据；独立校验器位于求解结果之后，
对外发布前必须通过复核。

\begin{figure}[htbp]
  \centering
  \begin{tikzpicture}[
    node distance=0.75cm,
    layer/.style={draw=black,fill=white,text=black,rounded corners=1pt,
      minimum width=12.4cm,minimum height=1.05cm,align=center,font=\small},
    item/.style={draw=black,fill=white,text=black,minimum height=0.72cm,
      align=center,font=\small,inner xsep=8pt},
    flow/.style={-{Latex[length=2mm]},draw=black,thick}
  ]
    \node[layer] (interaction) {应用交互层：决策驾驶舱 \quad 周课表 \quad 调课中心 \quad 资源管理};
    \node[layer,below=of interaction] (service) {应用服务层：REST API \quad 数据转换 \quad 方案评分 \quad 差异比较};
    \node[layer,below=of service] (decision) {决策核心层：CP-SAT 约束求解器 \quad 多策略目标 \quad 最小影响重排};
    \node[layer,below=of decision] (validation) {质量保障层：输入校验 \quad 独立课表校验 \quad 不可行诊断};
    \node[layer,below=of validation] (data) {业务数据层：教师 \quad 教室 \quad 班级 \quad 课程请求 \quad 时间槽};
    \node[item,left=0.8cm of service] (feishu) {飞书协同端\\（拟集成）};

    \draw[flow] (data) -- (validation);
    \draw[flow] (validation) -- (decision);
    \draw[flow] (decision) -- (service);
    \draw[flow] (service) -- (interaction);
    \draw[flow,<->] (feishu) -- (service);
  \end{tikzpicture}
  \caption{ClassMind 系统总体架构}
  \label{fig:architecture-new}
\end{figure}

\begin{table}[htbp]
  \centering
  \caption{系统模块及职责}
  \label{tab:modules-new}
  \small
  \begin{tabularx}{\textwidth}{L{2.6cm} L{3.1cm} Y}
    \toprule
    \textbf{模块} & \textbf{主要输入输出} & \textbf{职责} \\
    \midrule
    数据模型 & 业务对象、JSON 数据 & 统一字段定义，承载约束所需的结构化信息 \\
    约束求解器 & 问题实例、策略参数；课表 & 构造决策变量与约束，搜索并优化可行方案 \\
    独立校验器 & 问题实例、输出课表；错误清单 & 复核课次完整性、资源互斥、资质、容量和设备条件 \\
    评分与调课服务 & 课表、基线方案；评分和差异 & 计算统一指标，生成最小影响调课结果 \\
    API 与交互界面 & HTTP 请求；页面和 JSON & 连接算法与用户，展示方案、资源和变更明细 \\
    飞书适配层 & 多维表格、消息和审批事件 & 完成字段映射、方案回写和流程触发（第二阶段） \\
    \bottomrule
  \end{tabularx}
\end{table}

\subsection{约束模型}

设课次集合为 $L$，教师集合为 $T$，教室集合为 $R$，时间槽集合为 $S$。
定义布尔决策变量
\[
  x_{l,t,r,s}=\begin{cases}
  1, & \text{课次 $l$ 安排给教师 $t$、教室 $r$、时间槽 $s$；}\\
  0, & \text{否则。}
  \end{cases}
\]

每个课次必须且只能选择一个可行组合：
\[
  \sum_{t\in T}\sum_{r\in R}\sum_{s\in S}x_{l,t,r,s}=1,
  \qquad \forall l\in L.
\]

对任意时间槽，同一教师、教室和班级最多承担一个课次。例如教师互斥约束为
\[
  \sum_{l\in L}\sum_{r\in R}x_{l,t,r,s}\leq 1,
  \qquad \forall t\in T,\ s\in S.
\]
教师资质、教室容量、设备要求及各对象不可用时间在变量生成阶段预筛选；
不满足条件的组合不创建决策变量，从而缩小搜索空间。

软目标采用加权惩罚形式：
\[
  \min Z = w_p P_{\mathrm{pref}} + w_e P_{\mathrm{resource}}
  + w_b P_{\mathrm{balance}} + w_k P_{\mathrm{compact}}
  + w_c P_{\mathrm{change}}.
\]
不同策略只改变权重，不改变硬约束。跨策略比较时使用统一分项指标，而不直接
比较因权重不同而量纲不同的原始目标值。

\subsection{负载均衡与最小影响调课}

教师日负载均衡采用相对目标课次数的绝对偏差作为线性代理，而不是把它表述
为方差的一阶泰勒近似。设教师 $t$ 在日期 $d$ 的课次数为 $c_{t,d}$，目标
课次数为 $\bar c$，则引入非负辅助变量 $u_{t,d}$，并令
\[
  u_{t,d}\geq c_{t,d}-\bar c,\qquad
  u_{t,d}\geq \bar c-c_{t,d}.
\]
最小化 $\sum_{t,d}u_{t,d}$ 即可降低日负载的绝对偏差。

教师请假时，系统先读取原课表作为基线，再把新增不可用时段写入问题实例，
最后以变更课次数量为主要软目标重新求解。

\begin{algorithm}[htbp]
  \caption{最小影响调课流程}
  \label{alg:min-change-new}
  \begin{algorithmic}[1]
    \Require 原问题 $P$、基线课表 $B$、请假教师 $t^*$、时段 $s^*$
    \Ensure 新课表 $B'$、变更集合 $\Delta$
    \State 将 $s^*$ 加入 $t^*$ 的不可用时段，得到新问题 $P'$
    \State 以 $B$ 为基线，为每个发生变化的课次设置变更惩罚
    \State 求解 $P'$，并优先最小化变更课次数量
    \State 使用独立校验器复核新课表 $B'$
    \State 比较 $B$ 与 $B'$，生成变更集合 $\Delta$
    \State \Return $B',\Delta$
  \end{algorithmic}
\end{algorithm}

\subsection{飞书集成方案}

第二阶段拟以飞书多维表格作为业务数据维护入口：教师、教室、班级、课程和
时间槽分别映射为数据表；触发排课后，适配层读取记录并转换为统一 JSON；
求解通过后把课表和评分回写；教务确认方案后再触发审批和消息卡片。飞书
令牌、权限和回调验签统一在适配层处理，求解器不保存平台凭据。

\subsection{技术路线}

研究按“需求与数据字典—约束模型—原型实现—独立校验—实验评价—飞书集成”
推进。每一阶段均产生可验证交付物，模型和实现通过自动化测试连接，避免
界面演示与算法证据脱节。

% ==================== 第 5 章 重点难点、创新点与可行性 ====================
\clearpage
\section{重点难点、创新点与可行性}

\subsection{研究重点与难点}

\subsubsection{业务规则的形式化与一致性}

同一条业务规则可能同时出现在数据录入、求解模型、校验逻辑和页面提示中。
如果各层口径不一致，系统即使“求解成功”也可能无法执行。本课题的重点是
建立唯一的数据字典和约束清单，并让求解器与独立校验器基于同一业务定义、
采用相互独立的实现路径。

\subsubsection{多目标权重与方案解释}

多目标加权能够生成差异化方案，但权重并不天然具有业务含义。研究难点在于
保证指标量纲统一、权重变化能够反映策略意图，并避免用一个总分掩盖局部
退化。课题将保留学生偏好、教师偏好、资源适配、负载均衡和变更数量等分项
指标，并开展权重敏感性实验。

\subsubsection{规模扩展与动态重排}

排课搜索空间随课次、教师、教室和时间槽数量快速增长。课题将通过不可行
组合预筛选、对称性削减、求解时限和分区求解等方法控制规模。调课场景还要
在新增约束与方案稳定之间权衡，因此需要对“最少变更”建立清晰的度量和
验收方法。

\subsection{拟创新点}

\begin{enumerate}
  \item \textbf{求解与独立校验的双层保障。}求解器负责生成候选课表，
  校验器直接基于业务对象复核结果，使“可行”结论不依赖单一代码路径。
  \item \textbf{面向教务决策的多策略解释。}把不同角色关注点拆为统一分项
  指标，通过策略权重生成候选方案，并解释每个方案的得失，而不是只展示
  一个不可解释的总目标值。
  \item \textbf{基于基线差异的最小影响调课。}把原课表作为显式输入，以
  变更课次为惩罚项，输出新方案的同时生成逐条差异和受影响对象。
  \item \textbf{算法核心与协同平台解耦。}以稳定 API 封装排课能力，再接入
  多维表格、消息和审批，使求解逻辑可以独立测试，也便于替换交互端。
\end{enumerate}

上述创新点是面向具体应用的系统与方法创新，不声称对通用 CP-SAT 求解理论
作出突破。其有效性需要通过后续实验和用户评价验证。

\subsection{技术可行性}

OR-Tools CP-SAT 已提供布尔变量、整数变量、资源互斥和线性目标等建模能力，
能够直接表达本课题的核心约束。Python 适合快速实现数据模型、校验器和服务
接口；原生 Web 技术可满足原型展示；自动化测试则为持续修改约束提供回归
保障。当前可运行原型进一步证明了基础技术链路可打通。

\subsection{经济与实施可行性}

课题开发阶段以开源软件和本地运行环境为主，不依赖专用硬件。模块化设计使
算法、接口和前端可以逐步迭代；飞书集成放在核心求解稳定之后，可降低平台
权限、回调配置和接口变化对算法研发的干扰。实际部署成本将在真实规模实验
后，结合算力、运维和接口调用量重新评估。

\subsection{合规可行性}

演示数据使用虚构身份和资源信息。进入真实试点前，需要落实最小权限、访问
控制、传输加密、日志脱敏、数据保存期限和删除机制；飞书应用凭据不得进入
前端或源码仓库。若涉及未成年人信息，还应在组织授权和监护人知情范围内
处理，并仅采集完成排课所必需的字段。

% ==================== 第 6 章 前期基础与验证计划 ====================
\clearpage
\section{前期研究基础与验证计划}

\subsection{前期研究基础}

项目已形成第一阶段可运行原型。现有代码包含教师、教室、班级、课程请求和
时间槽数据模型；基于 CP-SAT 的排课求解器；独立课表校验器；候选策略与
评分服务；教师请假后的局部重排；本地 REST API；决策驾驶舱、周课表、
调课中心和资源概览四个页面。源码结构、接口样例和关键实现节选见附录。

原型采用一份小规模虚构数据进行功能验证，包含 4 位教师、3 间教室、4 个
班级、4 项课程请求和 8 个待排课次。该数据只用于验证模型和接口链路，不能
代表真实机构规模，也不据此推导行业效率或成本结论。

\subsection{当前可复核结果}

表 \ref{tab:verified-baseline} 仅列出可由当前仓库代码和自动化测试直接复核
的结果。求解耗时受设备、依赖版本和系统负载影响，因此不在此给出固定的
毫秒数；后续实验将同步记录运行环境和原始测量数据。

\begin{table}[htbp]
  \centering
  \caption{当前原型的可复核验证结果}
  \label{tab:verified-baseline}
  \small
  \begin{tabularx}{\textwidth}{L{3.0cm} L{3.2cm} Y}
    \toprule
    \textbf{验证对象} & \textbf{当前结果} & \textbf{复核依据} \\
    \midrule
    演示课次完整性 & 8 个课次均生成安排 & 求解结果课次数与课程请求课次总数一致 \\
    硬约束正确性 & 独立校验器报告 0 个冲突 & 教师、教室、班级互斥及资质、容量、设备检查 \\
    候选策略 & 学生、教师和资源效率三类方案可生成 & 方案接口及策略测试 \\
    局部重排 & 教师请假后返回新方案与差异 & 基线课表、请假约束和变更清单测试 \\
    自动化测试 & 19 项测试通过 & 12 项求解与服务测试、7 项 API 集成测试 \\
    页面与接口 & 4 个页面及主要 API 可访问 & 本地 HTTP 集成测试 \\
    \bottomrule
  \end{tabularx}
\end{table}

\subsection{实验设计}

后续实验分为五组：

\begin{enumerate}
  \item \textbf{正确性实验。}构造正常、边界和不可行实例，验证硬约束、
  数据错误提示和独立校验结果。
  \item \textbf{规模实验。}按课次数量、教师数量、教室数量和时间槽数量
  分层扩展数据，记录变量数、求解状态、首次可行解时间和总耗时。
  \item \textbf{策略敏感性实验。}逐步改变软目标权重，观察统一分项指标、
  课表结构和方案稳定性的变化，不比较不同量纲的原始目标值。
  \item \textbf{调课实验。}设置单教师请假、多时段不可用和教室停用事件，
  统计变更课次数、受影响班级和求解成功率。
  \item \textbf{可用性实验。}邀请目标用户完成数据导入、方案比较、调课和
  发布任务，记录任务完成率、操作时间和主观评价。
\end{enumerate}

\subsection{评价指标}

\begin{table}[htbp]
  \centering
  \caption{实验评价指标及计算口径}
  \label{tab:evaluation-metrics}
  \small
  \begin{tabularx}{\textwidth}{L{3.0cm} L{3.6cm} Y}
    \toprule
    \textbf{指标} & \textbf{计算口径} & \textbf{用途} \\
    \midrule
    硬冲突数 & 校验器发现的教师、教室、班级等冲突总数 & 判断方案是否可执行 \\
    课次完成率 & 已排课次除以请求课次总数 & 检查是否遗漏课程 \\
    偏好满足率 & 命中有效偏好时段的课次占比 & 衡量学生或教师体验 \\
    容量适配率 & 实际使用座位数除以分配教室容量 & 衡量资源匹配程度 \\
    负载绝对偏差 & 教师日课次相对目标值的绝对偏差之和 & 衡量日负载均衡 \\
    调课变更率 & 相对基线发生变化的课次占比 & 衡量局部重排稳定性 \\
    求解时间 & 在注明软硬件环境下的端到端耗时 & 评估规模扩展能力 \\
    任务完成率 & 用户在限定条件下完成指定操作的比例 & 评估系统可用性 \\
    \bottomrule
  \end{tabularx}
\end{table}

\subsection{结果记录与复现}

每次实验将保存输入数据、策略参数、依赖版本、随机种子、求解状态、指标和
原始输出。性能实验至少重复多次并报告中位数及波动范围；无法复现或缺少
来源的数据不进入最终结论。自动化测试在代码变更后持续运行，防止新增软目标
破坏已有硬约束。

% ==================== 第 7 章 进度、分工与预期成果 ====================
\clearpage
\section{进度安排、团队分工与预期成果}

\subsection{进度安排}

进度以“模型先稳定、集成后推进、实验贯穿全程”为原则。表
\ref{tab:research-plan} 中前期原型为已完成基础，其余为拟开展工作。

\begin{table}[htbp]
  \centering
  \caption{课题研究进度安排}
  \label{tab:research-plan}
  \small
  \begin{tabularx}{\textwidth}{L{2.3cm} L{2.8cm} Y Y}
    \toprule
    \textbf{阶段} & \textbf{时间} & \textbf{主要任务} & \textbf{阶段成果} \\
    \midrule
    前期基础 & 2026.07 上旬 & 完成数据模型、基础求解器、校验器、API 与页面原型 & 可运行原型和自动化测试 \\
    模型完善 & 2026.07--08 & 复核约束清单，完善权重、诊断和局部重排 & 约束说明、模型版本与回归测试 \\
    规模实验 & 2026.09 & 构造分层数据，开展性能与敏感性实验 & 原始实验数据和分析报告 \\
    飞书集成 & 2026.10 & 完成多维表格映射、方案回写和消息卡片原型 & 集成原型与接口文档 \\
    流程验证 & 2026.11 & 接入审批流程，开展调课与可用性测试 & 场景测试记录和问题清单 \\
    总结交付 & 2026.12 & 整理代码、文档、演示材料和研究结论 & 完整作品、报告和演示视频 \\
    \bottomrule
  \end{tabularx}
\end{table}

\subsection{团队分工}

当前报告按已确认负责人李帅锋编制，主要承担约束建模、求解器、服务接口、
前端原型、测试与文档整合。其他成员及指导教师信息应以正式报名表为准，
未确认姓名和贡献比例不在报告中设置占位数据。后续如增加成员，将按“任务、
产出、复核人”记录实际贡献，并在提交前统一更新封面和分工说明。

\subsection{预期成果与验收材料}

\begin{enumerate}
  \item 一套可运行的 ClassMind 智能排课原型及源代码；
  \item 一套覆盖硬约束、软目标、局部重排和独立校验的模型说明；
  \item 一组可复现实验数据、测试脚本和评价报告；
  \item 飞书多维表格、消息卡片和审批流程的集成原型与接口文档；
  \item 开题报告、结题报告、用户说明和答辩演示材料。
\end{enumerate}

\subsection{风险识别与应对}

图 \ref{fig:risk-matrix} 是项目组依据当前计划作出的定性风险评估：横轴为
发生概率，纵轴为影响程度。风险编号越靠近右上方，越需要优先跟踪。该图
使用白底、黑字和黑色边框，适合黑白打印。

\begin{figure}[htbp]
  \centering
  \begin{tikzpicture}[x=2.2cm,y=1.45cm,font=\small]
    % 3x3 matrix
    \draw[black,thick] (0,0) rectangle (3,3);
    \draw[black] (1,0)--(1,3) (2,0)--(2,3) (0,1)--(3,1) (0,2)--(3,2);
    \node[below] at (0.5,0) {低};
    \node[below] at (1.5,0) {中};
    \node[below] at (2.5,0) {高};
    \node[rotate=90,left] at (0,0.5) {低};
    \node[rotate=90,left] at (0,1.5) {中};
    \node[rotate=90,left] at (0,2.5) {高};
    \node[below=0.55cm] at (1.5,0) {发生概率};
    \node[rotate=90,left=0.75cm] at (0,1.5) {影响程度};

    \node[draw=black,fill=white,circle,inner sep=2pt] at (1.45,2.55) {R1};
    \node[draw=black,fill=white,circle,inner sep=2pt] at (1.75,2.20) {R2};
    \node[draw=black,fill=white,circle,inner sep=2pt] at (2.45,1.55) {R3};
    \node[draw=black,fill=white,circle,inner sep=2pt] at (1.35,1.35) {R4};
    \node[draw=black,fill=white,circle,inner sep=2pt] at (0.65,2.30) {R5};
  \end{tikzpicture}
  \caption{项目风险概率—影响矩阵}
  \label{fig:risk-matrix}
\end{figure}

\begin{table}[htbp]
  \centering
  \caption{主要风险及应对措施}
  \label{tab:risk-actions}
  \small
  \begin{tabularx}{\textwidth}{C{1.1cm} L{3.1cm} Y Y}
    \toprule
    \textbf{编号} & \textbf{风险} & \textbf{预警信号} & \textbf{应对措施} \\
    \midrule
    R1 & 真实规模下求解超时 & 首次可行解时间随规模快速增长 & 预筛选变量、设置时限、分区求解并保留可行解回退 \\
    R2 & 业务规则遗漏或冲突 & 试点数据频繁触发人工修正 & 建立约束清单，增加边界测试并由业务人员复核 \\
    R3 & 飞书接口或权限受限 & 回调、配额或字段权限不满足方案 & 保持核心服务独立，使用适配层和批处理降低耦合 \\
    R4 & 多目标权重难以解释 & 权重小幅变化导致方案大幅波动 & 使用统一分项指标，开展敏感性实验并提供默认策略 \\
    R5 & 数据隐私与合规风险 & 采集字段超出排课必要范围 & 最小化采集、脱敏演示、访问控制并设置数据生命周期 \\
    \bottomrule
  \end{tabularx}
\end{table}


% 参考文献
% ==================== 参考文献 ====================
\clearpage
\section*{参考文献}
\addcontentsline{toc}{section}{参考文献}

\begin{enumerate}[leftmargin=2em,itemsep=0.3em]
  \item Burke E K, Petrovic S. Recent research directions in automated
    timetabling[J]. European Journal of Operational Research,
    2002, 140(2): 266--280.

  \item de Werra D. Scheduling in sports[M]//Studies in Integer
    Programming. North-Holland, 1977: 381--395.

  \item Google OR-Tools Team. OR-Tools: Software for Combinatorial
    Optimization[EB/OL]. \url{https://developers.google.com/optimization},
    2024.

  \item Rossi F, van Beek P, Walsh T. Handbook of Constraint
    Programming[M]. Elsevier, 2006.

  \item Apt K. Principles of Constraint Programming[M]. Cambridge
    University Press, 2003.

  \item Laborie P, Rogerie J, Shaw P, et al. Reasoning with
    conditional time-intervals[C]//Proceedings of the Twenty-First
    International Florida Artificial Intelligence Research Society
    Conference. 2008: 555--560.

  \item 中国教育部. 关于进一步减轻义务教育阶段学生作业负担和
    校外培训负担的意见[Z]. 2021.

  \item 字节跳动飞书团队. 飞书开放平台开发者文档[EB/OL].
    \url{https://open.feishu.cn/document}, 2024.

  \item Cormen T H, Leiserson C E, Rivest R L, et al. Introduction
    to Algorithms (4th ed.)[M]. MIT Press, 2022.

  \item Taha H A. Operations Research: An Introduction (10th ed.)[M].
    Pearson, 2017.

  \item Python Software Foundation. The Python Language Reference
    (3.12)[EB/OL]. \url{https://docs.python.org/3.12/}, 2024.

  \item 维基百科. Constraint Programming[EB/OL].
    \url{https://en.wikipedia.org/wiki/Constraint_programming},
    2024.

  \item Russell S, Norvig P. Artificial Intelligence: A Modern
    Approach (4th ed.)[M]. Pearson, 2020.
\end{enumerate}


% 附录：大段代码、数据样例与接口细目统一置于正文之后
% ==================== 附录 ====================
\clearpage
\section*{附\quad 录}
\addcontentsline{toc}{section}{附\quad 录}

\subsection*{附录 A：核心数据类定义（节选）}
\addcontentsline{toc}{subsection}{附录 A：核心数据类定义（节选）}

本附录节选自 \code{classmind/models.py}，用于说明业务对象在代码层的
承载方式。完整源码请参见项目仓库 \code{classmind/models.py}。

\begin{lstlisting}[style=cmplain]
# classmind/models.py 节选
from __future__ import annotations
from dataclasses import asdict, dataclass, field
from typing import Any


@dataclass(frozen=True)
class TimeSlot:
    id: str
    day: str
    period: str
    order: int


@dataclass(frozen=True)
class Teacher:
    id: str
    name: str
    qualifications: tuple = ()
    unavailable: tuple = ()
    preferred_slots: tuple = ()


@dataclass(frozen=True)
class Room:
    id: str
    name: str
    campus: str
    capacity: int
    equipment: tuple = ()
    unavailable: tuple = ()


@dataclass(frozen=True)
class ClassGroup:
    id: str
    name: str
    size: int
    unavailable: tuple = ()
    preferred_slots: tuple = ()


@dataclass(frozen=True)
class CourseRequest:
    id: str
    name: str
    subject: str
    class_id: str
    sessions: int
    qualified_teacher_ids: tuple = ()
    required_equipment: tuple = ()
    allowed_slots: tuple = ()


@dataclass(frozen=True)
class ScheduledLesson:
    lesson_id: str
    course_id: str
    course_name: str
    class_id: str
    class_name: str
    teacher_id: str
    teacher_name: str
    room_id: str
    room_name: str
    slot_id: str
    day: str
    period: str
\end{lstlisting}

\subsection*{附录 B：演示数据示例（节选）}
\addcontentsline{toc}{subsection}{附录 B：演示数据示例（节选）}

以下片段来自 \code{data/demo.json}，仅保留字段结构与一两条记录，
完整数据请参见项目仓库 \code{data/demo.json}。

\begin{lstlisting}[style=cmjson]
// data/demo.json 节选
{
  "time_slots": [
    {"id":"MON_0800","day":"周一","period":"08:00-09:30","order":1},
    {"id":"FRI_1530","day":"周五","period":"15:30-17:00","order":20}
  ],
  "teachers": [
    {"id":"T01","name":"张老师","qualifications":["数学"],
     "unavailable":["WED_1000"],
     "preferred_slots":["MON_0800","TUE_0800","THU_1000"]},
    {"id":"T02","name":"李老师","qualifications":["数学","物理"],
     "unavailable":["TUE_1000"]}
  ],
  "rooms": [
    {"id":"R101","name":"思源教室","campus":"鼓楼校区",
     "capacity":30,"equipment":["投影","白板"]}
  ],
  "classes": [
    {"id":"C01","name":"高一数学A班","size":28}
  ],
  "courses": [
    {"id":"CR01","name":"高一数学提升","subject":"数学",
     "class_id":"C01","sessions":2,
     "qualified_teacher_ids":["T01","T02"],
     "required_equipment":["投影"]}
  ]
}
\end{lstlisting}

\subsection*{附录 C：API 接口详细定义}
\addcontentsline{toc}{subsection}{附录 C：API 接口详细定义}

本节列出当前原型已实现的核心接口，字段定义以服务实现为准。

\subsubsection*{C.1 GET /api/health}

\textbf{用途：}健康检查。\textbf{请求：}无。

\begin{lstlisting}[style=cmjson]
{
  "status": "ok",
  "service": "ClassMind",
  "solver": "OR-Tools CP-SAT"
}
\end{lstlisting}

\subsubsection*{C.2 GET /api/plans}

\textbf{用途：}获取多套候选方案。\textbf{请求：}无。
\textbf{响应：}包含方案列表，每条方案附 \code{status}、\code{metrics}、
\code{conflicts}、\code{schedule}、\code{scorecard} 字段。

\subsubsection*{C.3 GET /api/dashboard}

\textbf{用途：}获取决策驾驶舱所需数据。
\textbf{响应：}单个 \code{SolveResult}，额外含 \code{teacher\_load} 字段。

\subsubsection*{C.4 POST /api/solve}

\textbf{用途：}对自定义业务数据求解。
\textbf{请求体：}与 \code{data/demo.json} 同结构，可选 \code{strategy} 字段。
\textbf{响应：}单个 \code{SolveResult}。

\subsubsection*{C.5 POST /api/reschedule}

\textbf{用途：}教师请假最小影响调课。

\begin{lstlisting}[style=cmjson]
{ "teacher_id": "T01", "slot_id": "MON_0800" }
\end{lstlisting}

\textbf{响应：}包含 \code{before}、\code{after}、\code{diff}、\code{impact}
四个字段，分别表示调课前后方案、变更明细及影响对象。

\subsection*{附录 D：核心算法节选}
\addcontentsline{toc}{subsection}{附录 D：核心算法节选}

完整源码托管在 GitHub：
\begin{center}
  \code{github.com/Lenny-lab/classmind-scheduler}
\end{center}

本节给出 \code{solver.py} 求解入口与权重表片段，其余模块
（\code{api.py / service.py / validator.py / io.py / models.py / cli.py}）
请参见上述仓库的 \code{classmind/} 目录。

\begin{lstlisting}[style=cmplain]
# classmind/solver.py 节选
from ortools.sat.python import cp_model


STRATEGY_WEIGHTS = {
    "student":    {"preference": 10, "stability": 2,
                   "efficiency": 1, "change": 0, "compactness": 3},
    "teacher":    {"preference": 2,  "stability": 10,
                   "efficiency": 1, "change": 0, "compactness": 3},
    "efficiency": {"preference": 1,  "stability": 2,
                   "efficiency": 10,"change": 0, "compactness": 3},
    "balanced":   {"preference": 5,  "stability": 5,
                   "efficiency": 5, "change": 0, "compactness": 3},
    "reschedule": {"preference": 1,  "stability": 1,
                   "efficiency": 1, "change": 100, "compactness": 0},
}


def solve_schedule(problem, time_limit_seconds=15.0,
                   strategy="balanced", baseline=None):
    # 1. 业务数据校验
    data_errors = validate_problem(problem)
    if data_errors:
        return SolveResult(status="INVALID_DATA", conflicts=data_errors)

    # 2. 策略与权重选择
    if strategy not in STRATEGY_WEIGHTS:
        return SolveResult(status="INVALID_DATA",
                           conflicts=[f"未知策略: {strategy}"])
    weights = STRATEGY_WEIGHTS[strategy]

    # 3. 变量预筛选与硬约束建模（仅保留可行组合）
    model = cp_model.CpModel()
    variables = {}
    lesson_vars = defaultdict(list)
    for lesson_id, course in lesson_specs:
        for (teacher_id, room_id, slot_id) in feasible_combos(course):
            key = (lesson_id, teacher_id, room_id, slot_id)
            variables[key] = model.new_bool_var("x_" + "_".join(key))
            lesson_vars[lesson_id].append(variables[key])

    # 4. 互斥与"恰好一次"约束
    for lesson_id, _ in lesson_specs:
        model.add_exactly_one(lesson_vars[lesson_id])
    for resource in ("teacher", "room", "class"):
        for slot_id in slot_ids:
            model.add_at_most_one(...)

    # 5. 加权软目标
    objective_terms = []
    objective_terms += _per_lesson_soft_costs(...)
    objective_terms += _stability_cost(...)
    model.minimize(sum(objective_terms))

    # 6. 求解 + 独立校验
    solver = cp_model.CpSolver()
    solver.parameters.max_time_in_seconds = time_limit_seconds
    status_code = solver.solve(model)
    schedule = extract_schedule(solver, variables, ...)
    conflicts = validate_schedule(problem, schedule)
    return SolveResult(status=..., schedule=schedule,
                       conflicts=conflicts, metrics={...})
\end{lstlisting}

\subsection*{附录 E：术语表}
\addcontentsline{toc}{subsection}{附录 E：术语表}

\begin{table}[h]
  \centering
  \caption{术语表}
  \label{tab:glossary}
  \small
  \begin{tabularx}{\textwidth}{L{3.2cm} Y}
    \toprule
    \textbf{术语} & \textbf{说明} \\
    \midrule
    CP-SAT
      & Constraint Programming-SAT，Google OR-Tools 中的混合求解引擎，
        结合约束传播、SAT 冲突学习与线性规划松弛。\\
    CP & Constraint Programming，约束规划。\\
    CSP & Constraint Satisfaction Problem，约束满足问题。\\
    SAT & Boolean Satisfiability Problem，布尔可满足性问题。\\
    OR-Tools & Google 开源的运筹学工具包，Apache 2.0 协议。\\
    硬约束
      & 必须严格满足的约束，违反则方案无效（如教师同时段冲突数必须为 0）。\\
    软目标
      & 可以被违反但应尽量满足的目标（如学生偏好尽量命中）。\\
    策略权重
      & 不同维度的软目标在总目标中的相对重要性（0--10）。\\
    策略
      & 一组策略权重的命名（balanced / student / teacher / efficiency / reschedule）。\\
    稳定性绝对偏差
      & 教师日均课次相对目标值的绝对偏差之和，越低表示日负载越均匀。\\
    调课
      & 在原方案基础上做最小影响的修改。\\
    决策变量
      & 求解器可以控制的变量（如课次到时间槽的分配）。\\
    预筛选
      & 在生成决策变量前就剔除明显不可行的组合。\\
    辅助变量
      & 求解器引入的额外变量，用于表达复杂约束。\\
    教务老师
      & K12 培训机构中负责排课的老师。\\
    驾驶舱
      & Dashboard，集中展示关键指标的页面。\\
    \bottomrule
  \end{tabularx}
\end{table}


% 致谢
% ==================== 致谢 ====================
\clearpage
\section*{致\quad 谢}
\addcontentsline{toc}{section}{致\quad 谢}

本项目能够推进到此阶段，离不开多方面的支持与帮助，
在此一并致谢。

\begin{itemize}[leftmargin=2em,itemsep=0.2em]
  \item 感谢 Google OR-Tools 团队开源的 CP-SAT 求解器，
        这是本项目最核心的算法依赖；
  \item 感谢 Python 软件基金会维护的 Python 解释器与
        \code{unittest}、\code{http.server}、\code{dataclasses}、
        \code{json} 等标准库模块；
  \item 感谢 GitHub 提供免费的 Actions 持续集成服务，
        使本项目得以建立双版本自动测试流水线；
  \item 感谢飞书高校应用创新大赛组委会提供展示平台，
        明确了本项目后续集成与场景验证的方向；
  \item 感谢指导老师在选题、调研、方案设计各环节给予的悉心指导。
\end{itemize}

由于作者水平所限，文中难免存在不足之处，
敬请各位老师与读者批评指正。

\vspace{1em}
\begin{flushright}
  李帅锋\quad 敬上\\
  2026 年 7 月
\end{flushright}


\end{document}
